\documentclass[a4paper,12pt]{article}
\usepackage{docsTemplate}
\usepackage{longtable}

\setTitle{Verbale Interno 04/12/2025}
\setAuthors{Andrea Masiero}
\setVerificators{Bianca Zaghetto}
\setApprovation{Nenad Radulovic}
\setVersion{1.0}
\setType{Interno}
\setPlace{Discord}
\setTime{17.00}
\setDestination{Team Three Technologies}
\setAbsents{Sara Gioia Fichera, Francesco Balestro}

\begin{document}
    \include{memoTitlepage}

    \newpage
    \section*{Registro delle modifiche}

    \begin{table}[h!]
        \rowcolors{2}{lightgray}{white}
        \begin{tabularx}{\textwidth}{|l|l|l|l|X|X|}
            \hline
            \textbf{Vers.} & \textbf{Data} & \textbf{Autore} & \textbf{Ruolo} & \textbf{Verifica} & \textbf{Oggetto} \\
            \hline
            \textbf{1.0} & 09/12/2025 & Nenad Radulovic & Responsabile &  & Approvazione \\
            \hline
            \textbf{0.2} & 09/12/2025 & Andrea Masiero & Analista & Bianca Zaghetto & Correzioni \\
            \hline
            \textbf{0.1} & 06/12/2025 & Andrea Masiero & Analista & Bianca Zaghetto & Prima stesura \\
        \hline
        \end{tabularx}
    \end{table}

    \newpage
    \tableofcontents
    \newpage

    \section{Ordine del giorno} {
        \begin{itemize}
            \item Pianificazione scadenze
            \item PdP e PdQ
            \item Definizione dei macro
Casi d’Uso del sistema
        \end{itemize}
    }

    \section{Svolgimento} {
        \begin{longtable}{P{5cm}|P{10cm}} 
            \textbf{Pianificazione scadenze} & La discussione ha evidenziato un problema  nelle tempistiche: calcolando le 11 settimane previste dal 3 novembre, la consegna della milestone RTB cade intorno al 18 gennaio. Questa data si sovrappone alla sessione d'esami. Per gestire questo rischio, è stato stabilito come obiettivo primario il completamento dell'Analisi dei Requisiti entro dicembre.\\\\
            \textbf{PdP e PdQ} & Per quanto riguarda il Piano di Progetto (PdP), è stato deciso di impostarlo come uno strumento dinamico da aggiornare a ogni sprint. È stata assegnata la stesura di una bozza che dovrà focalizzarsi soprattutto sull'analisi dei rischi, come i ritardi legati alla sessione d'esami o l'apprendimento di nuove tecnologie.
Il Piano di Qualifica (PdQ) rappresenta invece l'aspetto più critico a causa delle scarse spiegazioni ricevute su test e metriche. Il gruppo ha deciso di sospenderne temporaneamente la lavorazione in attesa di chiarimenti da parte del docente.

\\\\
            \textbf{Definizione dei macro Casi d'Uso del sistema} & Sono stati definiti i macro Casi d'Uso, decidendo di unificare Visualizzazione e Caricamento e confermando funzionalità essenziali come Ricerca, Salvataggio e Verifica della firma. È stato stabilito di documentare formalmente gli errori (es. validazione) come deviazioni dai flussi principali, accantonando invece l'autenticazione.
Sul fronte interfaccia, si è proposta una visualizzazione\textbf{ }doppia (tramite checkbox) per alternare tra la struttura completa ad albero e una vista semplificata dei soli allegati. Per la stesura del documento, infine, si adotterà una numerazione gerarchica a tre livelli, puntando a un'elevata profondità di dettaglio ispirata ai migliori progetti degli anni passati.\\\\
        \end{longtable}
    }

    \section{Decisioni} {
        \rowcolors{2}{lightgray}{white}
        \begin{tabularx}{\textwidth}{|l|X|l|}
            \hline
            \textbf{Identificativo} & \textbf{Descrizione}\\
            \hline
            VI\_20251204.d1 & Pianificazione e scadenze \\
            \hline
            VI\_20251204d2 & Struttura e contenuto dell'Analisi dei Requisiti \\
            \hline
            VI\_20251204.d3 & Metodologia di lavoro e strumenti\\
            \hline
        \end{tabularx}
    }

    \section{Attività} {
        \rowcolors{2}{lightgray}{white}
        \begin{tabularx}{\textwidth}{|l|X|l|l|}
            \hline
            \textbf{Incaricato} & \textbf{Task Todo} & \textbf{Identificativo} & \textbf{Issue} \\
            \hline
            \textbf{Andrea Masiero} & Stesura Verbale Interno & VI\_20251204.a1 & \#3\\
            \hline
            \textbf{Tutti} & Aggiornamento Analisi dei Requisiti & VI\_20251204.a2 & \\
            \hline
            \textbf{Tutti} & Piano di Progetto & VI\_20251204.a3 & \\
            \hline
            \end{tabularx}
    }

\end{document}